% Part: incompleteness
% Chapter: arithmetization-syntax
% Section: coding-terms

\documentclass[../../../include/open-logic-section]{subfiles}

\begin{document}

\olfileid{inc}{art}{trm}
\olsection{पदांचे सांकेतीकरण}

\begin{explain}
पद हा चिन्हांच्या अनुक्रमाचा एक विशिष्ट प्रकार आहे: पदांसाठीच्या
रचनानियमांनुसार स्थिरचिन्हे आणि चरे यांपासून त्याची विगमनाधारित
रचना होते. चिन्हे सांकेतिक करण्याची पद्धत आणि संख्यांचे अनुक्रम
सांकेतिक करण्याची पद्धत वापरून चिन्हांचे अनुक्रम संख्यांनी सांकेतिक
करता येतात. त्यामुळे पदांना गोडेल संख्या देणे कठीण नाही. खरी
अडचण ही आहे की एखादी संख्या योग्य रीतीने रचलेल्या पदाची गोडेल
संख्या आहे हा गुणधर्म संगणनक्षम, किंबहुना आदिम पुनरावर्ती आहे, हे दाखवणे.
\end{explain}

!!^{variable}s आणि !!{constant}s ही सर्वांत सोपी पदे आहेत. $x$ ही
अशा पदाची गोडेल संख्या आहे का, हे तपासणे सोपे आहे:
$x$ ही एखाद्या~$i$ साठी $\Gn{\Obj v_i}$ असेल, तर आणि तरच
$\fn{Var}(x)$ खरे असते. दुसऱ्या शब्दांत, $x$ हा लांबी~$1$
असलेल्या अनुक्रमाचा संकेतांक आहे आणि त्यातील एकमेव घटक $(x)_0$
हा एखाद्या !!{variable}~$\Obj v_i$ चा संकेतांक आहे; म्हणजे एखाद्या~$i$
साठी $x$ हा $\tuple{\tuple{1, i}}$ आहे. तसेच, $x$ ही एखाद्या~$i$
साठी $\Gn{\Obj c_i}$ असेल, तर आणि तरच $\fn{Const}(x)$ खरे असते.
हे दोन्ही संबंध आदिम पुनरावर्ती आहेत; कारण असा~$i$ अस्तित्वात
असेल, तर तो $< x$ असलाच पाहिजे:
\begin{align*}
  \fn{Var}(x) & \defiff \bexists{i<x}{x = \tuple{\tuple{1, i}}}\\
  \fn{Const}(x) & \defiff \bexists{i<x}{x = \tuple{\tuple{2, i}}}
\end{align*}

\begin{prop}
\ollabel{prop:term-primrec}
$x$ ही अनुक्रमे पदाची किंवा बंद पदाची गोडेल संख्या असेल, तर आणि तरच
खरे असणारे $\fn{Term}(x)$ आणि $\fn{ClTerm}(x)$ हे संबंध आदिम पुनरावर्ती आहेत.
\end{prop}

\begin{proof}
चिन्हांचा अनुक्रम~$s$ हा पद असेल, तर आणि तरच पदांची अशी
रचनाक्रमिका~$s_0$, \dots, $s_{k-1} = s$ अस्तित्वात असते की जिच्यात
पदांच्या रचनानियमांनुसार !!{constant}s आणि !!{variable}s यांपासून
पद~$s$ कसे बनले याची नोंद असते. एखादी प्रस्तावित रचनाक्रमिका
रचनानियमांचे पालन करते ही अट व्यक्त करण्यासाठी, प्रत्येक $i < k$
बाबत पुढीलपैकी एक अट खरी असली पाहिजे:
\begin{enumerate}
\item $s_i$ हे !!a{variable}~$\Obj v_j$ आहे; किंवा
\item $s_i$ हे !!a{constant}~$\Obj c_j$ आहे; किंवा
\item $i$-व्या स्थानापूर्वी आलेल्या $t_1$, \dots, $t_n$ या $n$
  पदांपासून $n$-स्थानी !!{function}~$\Obj f^n_j$ वापरून $s_i$ बनले आहे.
\end{enumerate}
गोडेल संख्यांवरील संबंधित संबंध आदिम पुनरावर्ती आहे हे दाखवण्यासाठी
ही अट आदिम पुनरावर्ती फलने, संबंध आणि परिबद्ध संख्यकीकरण वापरून
व्यक्त करावी लागेल.

$y$ ही $s_0$, \dots, $s_{k-1}$ या अनुक्रमाचा संकेतांक आहे असे समजू;
म्हणजे $y = \tuple{\Gn{s_0}, \dots, \Gn{s_{k-1}}}$. तो गोडेल
संख्या~$x$ असलेल्या पदाच्या रचनाक्रमिकेचा संकेतांक असेल, तर आणि तरच
प्रत्येक $i < k$ साठी पुढीलपैकी एक अट खरी असते:
\begin{enumerate}
\item $\fn{Var}((y)_i)$; किंवा
\item $\fn{Const}((y)_i)$; किंवा
\item एखादा~$n$ आणि $z = \tuple{z_1, \dots, z_n}$ ही संख्या
  अस्तित्वात असतात, अशी की प्रत्येक $z_l$ हा $i' < i$ असलेल्या
  एखाद्या निर्देशांकासाठी $(y)_{i'}$ बरोबर असतो आणि
\[
(y)_i = \Gn{\Obj f^n_j(} \concat \fn{flatten}(z) \concat \Gn{)},
\]
\end{enumerate}
आणि याशिवाय $(y)_{k-1} = x$ असते. ($\fn{flatten}(z)$ हे फलन
$\tuple{\Gn{t_1}, \dots, \Gn{t_n}}$ या अनुक्रमापासून
$\Gn{t_1, \dots, t_n}$ तयार करते आणि ते आदिम पुनरावर्ती आहे.)

(3) मधील $j$, $n$ हे निर्देशांक, $t_l$ या पदांच्या गोडेल संख्या
$z_l$, आणि $\tuple{z_1, \dots, z_n}$ या अनुक्रमाचा संकेतांक~$z$
हे सर्व~$y$ पेक्षा लहान आहेत. वरच्या~$k$ ऐवजी $\len{y}$
वापरता येते. म्हणून ``$y$ ही गोडेल संख्या~$x$ असलेल्या पदाच्या
रचनाक्रमिकेचा संकेतांक आहे'' हे वाक्य त्या संबंधाचे आदिम पुनरावर्ती
असणे स्पष्ट करणाऱ्या रीतीने व्यक्त करता येते.

आता~$y$ साठी आदिम पुनरावर्ती मर्यादा आहे याचीच खात्री करून घ्यायची
आहे. $x$ ही एखाद्या पदाची गोडेल संख्या असेल, तर त्याला जास्तीत
जास्त $\len{x}$ पदांची रचनाक्रमिका असली पाहिजे: कारण~$s$ च्या
केवळ उपपदांची रचनाक्रमिका निवडता येते. तिच्यातील प्रत्येक पदाची
सुरुवात~$s$ मधील एखाद्या स्थानावर
होते, आणि दोन उपपदे एकाच स्थानावर सुरू होऊ शकत नाहीत. $s$ च्या
प्रत्येक उपपदाची गोडेल संख्या अर्थातच $\le x$ आहे. म्हणून
$k=\len{x}$ असेल, तर संकेतांक $\le p_{k-1}^{k(x+1)}$ असलेली
रचनाक्रमिका नेहमी अस्तित्वात असते.
%READERNOTE{SOL6-A063 — मर्यादा सिद्ध करताना गरज नसलेली पदे वगळून केवळ अंतिम पदाच्या उपपदांची रचनाक्रमिका निवडली आहे. मनमानी रचनाक्रमिकेत इतर मोठी पदे असू शकतात; त्यांच्यासाठी ही लांबीची मर्यादा दावा केलेली नाही.}

$\fn{ClTerm}$ साठी !!{variable}s संबंधी अट फक्त वगळावी.
\end{proof}

\begin{prob}
$\tuple{\Gn{t_1}, \dots, \Gn{t_n}}$ या अनुक्रमापासून
$\Gn{t_1, \dots, t_n}$ तयार करणारे $\fn{flatten}(z)$ हे फलन
आदिम पुनरावर्ती आहे, हे दाखवा.
\end{prob}

\begin{prop}\ollabel{prop:num-primrec}
  $\fn{num}(n) = \Gn{\num{n}}$ हे फलन आदिम पुनरावर्ती आहे.
\end{prop}

\begin{proof}
  $\fn{num}(n)$ ची आदिम पुनरावर्तनाने व्याख्या करू:
  \begin{align*}
    \fn{num}(0) & = \Gn{\Obj 0}\\
    \fn{num}(n+1) & = \Gn{\prime(} \concat \fn{num}(n) \concat \Gn{)}.
  \end{align*}
\end{proof}

\end{document}
